\section{Experiments}\label{sec:experiments}

We evaluate whether adapter-residual credit assignment gives a useful
token-level signal in the low-rank fine-tuning regime. The empirical section
focuses on a single controlled sweep, summarized in
Table~\ref{tab:sweep}. It uses one model, one task, one advantage estimator,
and seven weighting configurations. This design trades breadth for a cleaner
test of the central claim: ARCA changes credit assignment without adding
trainable parameters beyond the LoRA adapter used by the matched baselines.

\subsection{Setup}

\paragraph{Model and task.} We fine-tune \texttt{Qwen/Qwen3-1.7B-Base} on
the MATH training split and evaluate on held-out MATH examples. Prompts ask
the model to solve the problem step by step and place the final answer in
\texttt{\textbackslash boxed\{\}}. Rewards are binary exact-match scores
computed from the final boxed answer after normalization.

\paragraph{Training.} All runs use GRPO prompt-level advantages, LoRA
adapters on the attention and MLP projections, and the same rollout budget,
optimizer, learning-rate schedule, prompt format, and decoding parameters.
Each training step samples $K=4$ completions per prompt, uses 100 update
steps, and evaluates greedy accuracy and pass@4 on 200 held-out MATH
examples. We use a single seed (1337). We save checkpoints every 20 steps.
We use AdamW with learning rate $5\times10^{-6}$, zero weight decay, a
10-step warmup followed by cosine decay, batch size 2, gradient accumulation
over 4 microbatches, maximum prompt length 512, maximum generation length
512, temperature 1.0, and top-$p$ 0.95 sampling during training.

\paragraph{Methods.} The sweep contains seven runs. Four baselines use LoRA
rank $r=64$: uniform token weighting, surprisal weighting,
entropy-reduction weighting, and policy-divergence weighting. The proposed
method, ARCA, is run at $r \in \{4,16,64\}$. The $r=64$ comparison tests
ARCA against baselines with the same trainable parameter count. The $r=4$
and $r=16$ comparisons test whether ARCA remains competitive with fewer
adapter parameters than the $r=64$ baselines.

\begin{table*}[t]
\centering
\caption{\textbf{Reported sweep.} All runs use
\texttt{Qwen/Qwen3-1.7B-Base}, MATH, GRPO, one seed, and 100 training steps.
The non-ARCA baselines are matched at LoRA rank 64; ARCA is evaluated at
three ranks to separate credit assignment from adapter capacity.}
\label{tab:sweep}
\begin{tabular}{l c c c}
\toprule
Method & LoRA rank & Trainable params & Role in comparison \\
\midrule
Uniform & 64 & 69.7M & Same-rank baseline \\
Surprisal & 64 & 69.7M & Intrinsic token baseline \\
Entropy reduction & 64 & 69.7M & Intrinsic token baseline \\
Policy divergence & 64 & 69.7M & Intrinsic token baseline \\
ARCA & 4 & 4.36M & Low-rank ARCA control \\
ARCA & 16 & 17.4M & Mid-rank ARCA control \\
ARCA & 64 & 69.7M & Same-rank ARCA comparison \\
\bottomrule
\end{tabular}
\end{table*}

\subsection{Main Results on MATH}

Table~\ref{tab:math-results} reports held-out MATH performance after RL
fine-tuning. The primary comparison is ARCA at rank 64 versus the rank-64
baselines, which holds the trainable parameter count fixed. The rank-4 and
rank-16 ARCA runs provide a more stringent control: if lower-rank ARCA is
competitive with rank-64 baselines, the gain cannot be attributed simply to
using a larger adapter.

\begin{table*}[t]
\centering
\caption{\textbf{Held-out MATH performance.} Greedy accuracy and pass@4
after GRPO fine-tuning. ARCA at rank 64 is parameter-matched to the
rank-64 baselines; lower-rank ARCA rows test the parameter-count
alternative. Train reward is averaged over the 100 update steps.}
\label{tab:math-results}
\begin{tabular}{l c c c c}
\toprule
Method & LoRA rank & Greedy acc. & Pass@4 & Avg. train reward \\
\midrule
Uniform & 64 & \textbf{0.620} & 0.675 & 0.362 \\
Surprisal & 64 & 0.520 & 0.650 & 0.326 \\
Entropy reduction & 64 & 0.600 & \textbf{0.700} & 0.367 \\
Policy divergence & 64 & 0.565 & 0.670 & 0.351 \\
ARCA & 4 & 0.525 & 0.640 & 0.327 \\
ARCA & 16 & 0.510 & 0.655 & 0.336 \\
ARCA & 64 & 0.590 & 0.680 & 0.353 \\
\bottomrule
\end{tabular}
\end{table*}

Table~\ref{tab:math-results} should be read as a compact validation rather
than a broad benchmark. Uniform weighting obtains the highest greedy
accuracy, while entropy reduction obtains the highest pass@4. ARCA at rank
64 is competitive with these baselines under the same trainable parameter
count, trailing uniform greedy accuracy by three points and exceeding uniform
pass@4 by 0.5 points. The lower-rank ARCA variants do not outperform the
rank-64 baselines in this run, which is consistent with a capacity-performance
tradeoff rather than evidence that ARCA gains come from extra parameters.

\subsection{Credit-Assignment Diagnostics}

Performance alone does not show whether a method changes token-level credit
assignment. We therefore log the concentration of token weights during
training. We report the Gini coefficient $G(w)$ and effective-token ratio
$\mathrm{EffN}(w)/T$. Uniform weighting has $G=0$ and
$\mathrm{EffN}/T=1$; highly concentrated weights have $G$ near one and a
small effective-token ratio.

\begin{table}[t]
\centering
\caption{\textbf{Token-weight diagnostics.} Concentration statistics
averaged over training. These diagnostics test whether a method produces a
non-uniform but non-degenerate token-credit signal, the empirical observable
predicted by the degeneration analysis.}
\label{tab:weight-diagnostics}
\begin{tabular}{l c c c}
\toprule
Method & LoRA rank & $G(w)$ & $\mathrm{EffN}/T$ \\
\midrule
Uniform & 64 & 0.000 & 1.000 \\
Surprisal & 64 & 0.941 & 0.052 \\
Entropy reduction & 64 & 0.920 & 0.073 \\
Policy divergence & 64 & 0.939 & 0.055 \\
ARCA & 4 & 0.353 & 0.467 \\
ARCA & 16 & 0.388 & 0.460 \\
ARCA & 64 & 0.490 & 0.334 \\
\bottomrule
\end{tabular}
\end{table}

Table~\ref{tab:weight-diagnostics} is the central diagnostic result. Uniform
weighting is exactly flat. Surprisal, entropy reduction, and policy divergence
produce extremely concentrated distributions, with effective-token ratios
between 0.052 and 0.073 on average. These methods therefore do not collapse to
uniform in this implementation; they degenerate in the other direction, into
spuriously sparse credit. ARCA occupies the predicted middle regime: it is far
from uniform, but it does not concentrate credit onto the tiny effective
support used by the output-distribution baselines. This is the empirical
signature of the theoretical distinction. Output-distribution scores ask
where the model is uncertain or shifted; ARCA asks where the adapter actually
acts.

\subsection{Results and Discussion}

The results support a theory-first interpretation. The downstream accuracy
numbers do not show a universal ARCA win: uniform is best under greedy
decoding, and entropy reduction is best under pass@4. However, ARCA at rank
64 is competitive under matched parameters, and its token-credit distribution
is qualitatively different from every baseline. In this paper, that diagnostic
is not secondary. It is the observable predicted by the LoRA-degeneration
analysis: a method whose salience is measured through output distributions
can broadcast credit uniformly or concentrate it on a very small set of
positions, while adapter-residual salience remains position-discriminative
without becoming maximally sparse.

This distinction matters because token-level credit assignment is often
evaluated only by final task accuracy. The concentration diagnostics separate these cases.
ARCA demonstrates that one can obtain an adaptation-aware credit signal from
the LoRA adapter itself, with no value head, learned process reward model, or tree
construction, while keeping the trainable parameter count identical to the
rank-64 baselines.

The rank controls provide a useful boundary on the claim. ARCA at ranks 4 and
16 uses substantially fewer trainable parameters than the rank-64 baselines,
and in this sweep those lower-rank variants do not match rank-64 performance.
We therefore do not claim that ARCA removes the need for sufficient adapter
capacity. The supported claim is sharper: at a fixed adapter rank, ARCA
changes the geometry of token credit in the way predicted by the theory, and
does so while remaining competitive on held-out MATH.

\subsection{Implementation and Metrics}

For each minibatch, training samples completions from the current policy,
computes binary rewards after full generation, constructs GRPO advantages,
and applies the weighted token-level objective from
Section~\ref{sec:methods}. Surprisal uses current-policy token log
probabilities, entropy-reduction uses next-token entropy drops,
policy-divergence compares adapted and adapter-disabled log probabilities,
and ARCA weights tokens by adapter-residual norms in the final hidden layer.

The code records per-step reward mean, reward variance, loss, completion
length, token-weight Gini, effective-token ratio, and, for ARCA, the mean
and maximum adapter-residual norm. Final evaluation records greedy accuracy
and pass@4 on held-out MATH examples. These metrics allow us to distinguish
three claims: whether the method trained, whether it changed token-credit
structure, and whether that structure improved held-out performance.

\subsection{Compute}

Each run fits on a single H100-class GPU. The seven submission-time runs are
executed independently, one configuration per GPU, with wall-clock training
time on the order of a few hours per run plus final evaluation. The reported
sweep therefore requires seven single-GPU runs; preliminary development runs
were used to tune runtime and checkpointing but are not included in the
reported comparison.

\subsection{Broader Impacts}

This work is methodological and studies how to assign token-level credit
during reinforcement learning for language models. The potential positive
impact is improved sample efficiency and interpretability of RL fine-tuning,
especially in settings where practitioners already use parameter-efficient
adaptation. The same techniques could also improve the fine-tuning of models
for harmful or misleading generation if applied without appropriate task,
data, and deployment safeguards. We do not release a new general-purpose
model in this paper.

\subsection{Limitations}

This sweep is deliberately small: one model, one dataset, one seed, one
advantage estimator, and seven configurations. It is sufficient to test the
parameter-matched ARCA comparison and the low-rank control, but the reported
numbers are descriptive rather than statistical significance claims. We view
the results as a targeted validation of the mechanism rather than a broad
benchmark claim.
